\clearpage\SourceLeaf{177}%
\pdfbookmark[1]{Tafel für die Sonnen-Coordinaten}{sonnen1812}
\fancyhead[L]{\small TAFEL FÜR DIE SONNEN-COORDINATEN}
\begin{center}
{\small TAFEL FÜR DIE SONNEN-COORDINATEN.\par}
\medskip\hrule\medskip
{\small Monatliche Correspondenz zur Beförderung der Erd- und Himmels-Kunde,\par
herausgegeben vom Freiherrn \textsc{von Zach}. Januar 1812.\par}
\medskip\hrule
\end{center}\medskip
Die von mir zuerst im Jahre 1804 vorgeschlagene und nachher in der \emph{Theoria motus corporum coelestium} vervollkommnete Methode, die geocentrischen Örter der Himmelskörper durch ihre \emph{unmittelbare} Beziehung auf den Äquator, vermittelst dreier rechtwinklichter Coordinaten zu berechnen, erleichtert die Arbeit in allen Fällen, wo \emph{viele} Örter zu berechnen sind, in einem so hohen Grade, dass sie von den sachkundigen Astronomen mit Beifall aufgenommen und bereits vielfach in Gebrauch gekommen ist. Die Fälle, wo \emph{viele} geocentrische Örter zu berechnen vorkommen, sind zwiefach: entweder will man viele während einer mässigen Zwischenzeit gemachte Beobachtungen eines Himmelskörpers mit den Elementen vergleichen, um daraus den mittlern Fehler der letztern mit aller möglichen Schärfe zu bestimmen, oder man will die geocentrische Bewegung eines Himmelskörpers zur Erleichterung der Beobachtungen oder zu andern Zwecken voraus berechnen. In dem letztern Fall bedarf es einer sehr grossen Schärfe nicht, man will nicht --- kann vielleicht auch nicht --- Secunden verbürgen, sondern begnügt sich mit ganzen Minuten, und hier ist es, wo die Methode durch eine Hülfstafel noch einer sehr bedeutenden Erleichterung fähig ist. Wenn man nemlich, was \emph{hier} erlaubt ist, die Störungen der Länge der Sonne und des Abstandes von der Erde vernachlässigt, und zuerst für Excentricität der Erdbahn, Länge des Perigäum und Schiefe der Ecliptik bestimmte Werthe zum Grunde
\SourceLeaf{178}%
legt, so sind die drei Coordinaten der Sonne blos von der mittlern Sonnenlänge abhängig, und lassen sich daher bequem in eine Tafel bringen. Eine solche Tafel habe ich zu meinem eigenen Gebrauche bereits vor mehreren Jahren berechnet, indem ich die Elemente der Sonnenbahn für 1800 zum Grunde legte und der Bequemlichkeit wegen die mittlere Sonnenlänge durch einzelne Decimal-Grade wachsen liess. Es brauchte so nur noch eine Tafel für die Epochen des Arguments zu Anfang jedes Jahres und dessen Änderung für einzelne Tage u.\ s.\ w.\ hinzuzukommen. So war die Tafel auf eine Reihe von Jahren hinreichend genau, so lange nemlich die Argumente der Sonnenbahn von denen von 1800 noch nicht zu sehr verschieden waren. Um die Tafel auch für entlegenere Zeiten brauchen zu können, berechnete ich noch eine Tafel für die Säcular-Änderungen der Coordinaten, welche durch die Säcularänderungen der Elemente der Sonnenbahn hervorgebracht werden. Auf diese Weise konnte meine Tafel mehrere Jahrhunderte hindurch dienen.

So wie sich die einfachsten Methoden nicht immer zuerst darbieten, so war es auch hier gewesen. Hätte ich die Argumente der Tafel nicht um Einen Decimal-Grad, sondern um die mittlere tägliche Bewegung der Sonne wachsen lassen, so würde der Gebrauch der Tafel dadurch noch bedeutend erleichtert worden sein. Die \emph{unmittelbar} in der Tafel gegebenen Coordinaten hätten dann (abgesehen von den Säcular-Veränderungen) in jedem bestimmten Jahre jede für ein bestimmtes Datum und \emph{alle für einerlei Stunde} gegolten, und dadurch würden die vorläufigen Rechnungen zur Bildung der Argumente so gut wie ganz weggefallen sein. Ich veranlasste daher den geschickten Herrn \textsc{Wachter}, die Tafel nach diesem Princip ganz neu zu berechnen, und rieth ausserdem, da das Jahr 1800 sich nun immer mehr von uns entfernt, lieber die Elemente für das Jahr 1820 zum Grunde zu legen, damit man vorerst eine Reihe von Jahren hindurch (etwa von 1810 bis 1830) gar nicht nöthig hätte, die Tafel für die Säcular-Veränderungen mit zu berücksichtigen. Die Argumente der Tafel schreiten also immer mit Incrementen von $59'\ 8''33$ fort, sind aber so gewählt (der bequemern Berechnung wegen), dass sie nicht von 0 anfangen, sondern so dass die Länge des Perihels selbst ein Glied derselben wird. Diese von Hrn \textsc{Wachter} mit grosser Sorgfalt berechnete Tafel mache ich hier bekannt, und glaube dadurch den rechnenden Astronomen einen Dienst zu erweisen. Es bleibt mir nur übrig, einige Worte über den Gebrauch derselben hinzuzusetzen.
\SourceLeaf{179}%

Die I.\ Tafel gibt die Epoche, d.\ i.\ die Stunden u.\ s.\ w.\ nach Pariser Meridian, für welche die Tafel II.\ bei dem als ihr Argument beigefügten Datum gilt. Man muss also die Epochen der ersten Tafel zu $0^{\mathrm h}$ des Datums hinzusetzen, (blos in den Jahren 1792, 1793, 1796, 1797 verwandelt sich die Addition in Subtraction), um die Zeitpunkte zu haben, wo die Tafel II.\ unmittelbar gilt. Bei Berechnung von Ephemeriden für Planeten und Cometen, wo es ziemlich gleichgiltig ist, welche Stunden man wählt, kann man es also immer so einrichten, dass man die Tafel II.\ unmittelbar ohne Interpolation gebrauchen kann. Oder man kann auch, um wenigstens eine \emph{bequeme} Interpolation zu haben, solche Zeiten wählen, die von denen, für welche die Tafel unmittelbar gilt, um $\frac12$, $\frac13$ oder $\frac14$ Tag abliegen.

Will man hingegen die Sonnen-Coordinaten für eine vorgeschriebene Stunde haben, so zieht man zuerst das, was die Epochentafel gibt, ab (welches also in den Jahren 1792, 1793, 1796, 1797 eine Addition wird) und geht mit dem Rest in die Tafel II.\ ein.

\emph{Beispiel:} Man verlangt die Sonnen-Coordinaten für 1811 April 11 $7^{\mathrm h}47^{\mathrm m}18^{\mathrm s}$ M.\ Z.\ in Paris (durch Aberration verbesserte Beobachtungszeit für die erste Beobachtung des Cometen von Herrn \textsc{von Zach}). Man hat also
\begin{center}
\begin{tabular}{lr}
& April $11.\ 7^{\mathrm h}\ 47^{\mathrm m}18^{\mathrm s}=$ April $11.3245$\\
Epoche für 1811 & $.\quad.\quad.\quad1\quad5\quad7\quad26\quad.\quad.\quad.\quad.\quad.\quad.\quad1.2135$\\\cline{2-2}
& April $10\quad2\quad39\quad52=$ April $10.1110$
\end{tabular}
\end{center}
Damit gibt Herrn \textsc{Wachters} Tafel
\[
x=+9362.5,\qquad y=+3299.4,\qquad z=+1432.1
\]

Bei der III.\ Tafel, für die Säcular-Veränderung der Coordinaten, ist zu bemerken, dass ihre Zeichen nicht absolut, sondern algebraisch zu verstehen sind, also für ein Jahr \emph{nach} 1820 ein negatives $x$ durch die Säcular-Veränderung eine absolute Vergrösserung erleidet. Will man darauf in unserm Beispiele Rücksicht nehmen, so muss noch die Säcular-Veränderung mit $-\frac9{100}$ multiplicirt hinzugefügt werden, also
\[
\mathrm dx=+0.5,\qquad\mathrm dy=-0.1,\qquad\mathrm dz=-0.1
\]
folglich
\[
x=+9363.0,\qquad y=+3299.3,\qquad z=+1432.0
\]
\SourceLeaf{180}%

Übrigens gibt die Tafel alles für die mittlere Distanz der Erde von der Sonne $=10000$, also nach gewöhnlicher Schreibart
\[
x=+0.93630,\qquad y=+0.32993,\qquad z=+0.14320
\]
Herr \textsc{von Lindenau} hat berechnet
\[
x=+0.936308,\qquad y=+0.329822,\qquad z=+0.143148
\]

Offenbar kann man für mehr als vier Decimalen nicht einstehen, und es wäre überflüssig gewesen, die Tafel II.\ mit mehrern Decimalen zu geben, da die periodischen Störungen sogar die vierte Decimale um eine oder zwei Einheiten afficiren können. Allein \emph{zu dem Zweck}, wozu die Tafel bestimmt ist, ist diese Genauigkeit auch vollkommen hinreichend, und für die Fälle, wo die Genauigkeit von Secunden verlangt wird, kann und soll sie nicht gelten.
\clearpage
\SourceLeaf{181}%
\begin{center}
\emph{Tafel für die Coordinaten der Sonne.}

\medskip
I. Tafel. Die Epochen.

\medskip
\begingroup\fontsize{9}{10.2}\selectfont
\setlength{\tabcolsep}{3pt}
\begin{tabular}{r|rrrr|r||r|rrrr|r}
&\multicolumn{5}{c||}{Mittlere Pariser Zeit.}& &\multicolumn{5}{c}{Mittlere Pariser Zeit.}\\\hline
$1790$ & $+0^{\mathrm d}$ & $3^{\mathrm h}$ & $1^{\mathrm m}$ & $22^{\mathrm s}$ & $+0^{\mathrm d}1260$ &  &  &  &  &  & \\
$1791$ & $+0$ & $8$ & $50$ & $14$ & $+0.3682$ & $1831$ & $+1^{\mathrm d}$ & $1^{\mathrm h}$ & $24^{\mathrm m}$ & $37^{\mathrm s}$ & $+1^{\mathrm d}0588$\\
$1792$ & $-0$ & $9$ & $20$ & $54$ & $-0.3895$ & $1832$ & $0$ & $7$ & $13$ & $29$ & $0.3010$\\
$1793$ & $-0$ & $3$ & $32$ & $3$ & $-0.1473$ & $1833$ & $0$ & $13$ & $2$ & $20$ & $0.5433$\\
$1794$ & $+0$ & $2$ & $16$ & $49$ & $+0.0950$ & $1834$ & $0$ & $18$ & $51$ & $12$ & $0.7855$\\
$1795$ & $+0$ & $8$ & $5$ & $40$ & $-0.3373$ & $1835$ & $1$ & $0$ & $40$ & $3$ & $1.0278$\\[3pt]
$1796$ & $-0$ & $10$ & $5$ & $28$ & $-0.4205$ & $1836$ & $0$ & $6$ & $28$ & $55$ & $0.2701$\\
$1797$ & $-0$ & $4$ & $16$ & $37$ & $-0.1782$ & $1837$ &  & $12$ & $17$ & $47$ & $0.5123$\\
$1798$ & $+0$ & $1$ & $32$ & $15$ & $+0.0640$ & $1838$ &  & $18$ & $6$ & $38$ & $0.7546$\\
$1799$ &  & $7$ & $21$ & $7$ & $0.3063$ & $1839$ &  & $23$ & $55$ & $30$ & $0.9969$\\
$1800$ &  & $13$ & $9$ & $58$ & $0.5486$ & $1840$ &  & $5$ & $44$ & $21$ & $0.2391$\\[3pt]
$1801$ & $0$ & $18$ & $58$ & $50$ & $0.7908$ & $1841$ &  & $11$ & $33$ & $13$ & $0.4814$\\
$1802$ & $1$ & $0$ & $47$ & $42$ & $1.0331$ & $1842$ &  & $17$ & $22$ & $5$ & $0.7237$\\
$1803$ & $1$ & $6$ & $36$ & $33$ & $1.2754$ & $1843$ &  & $23$ & $10$ & $56$ & $0.9659$\\
$1804$ & $0$ & $12$ & $25$ & $24$ & $0.5176$ & $1844$ &  & $4$ & $59$ & $48$ & $0.2082$\\
$1805$ & $0$ & $18$ & $14$ & $16$ & $0.7599$ & $1845$ &  & $10$ & $48$ & $39$ & $0.4505$\\[3pt]
$1806$ & $1$ & $0$ & $3$ & $8$ & $1.0022$ & $1846$ &  & $16$ & $37$ & $31$ & $0.6927$\\
$1807$ & $1$ & $5$ & $51$ & $59$ & $1.2444$ & $1847$ &  & $22$ & $26$ & $22$ & $0.9350$\\
$1808$ & $0$ & $11$ & $40$ & $51$ & $0.4867$ & $1848$ &  & $4$ & $15$ & $14$ & $0.1773$\\
$1809$ & $0$ & $17$ & $29$ & $42$ & $0.7290$ & $1849$ &  & $10$ & $4$ & $6$ & $0.4195$\\
$1810$ & $0$ & $23$ & $18$ & $34$ & $0.9712$ & $1850$ &  & $15$ & $52$ & $57$ & $0.6618$\\[3pt]
$1811$ & $1$ & $5$ & $7$ & $26$ & $1.2135$ & $1851$ &  & $21$ & $41$ & $49$ & $0.9041$\\
$1812$ & $0$ & $10$ & $56$ & $17$ & $0.4557$ & $1852$ &  & $3$ & $30$ & $40$ & $0.1463$\\
$1813$ & $0$ & $16$ & $45$ & $9$ & $0.6980$ & $1853$ &  & $9$ & $19$ & $32$ & $0.3886$\\
$1814$ & $0$ & $22$ & $34$ & $0$ & $0.9403$ & $1854$ &  & $15$ & $8$ & $23$ & $0.6308$\\
$1815$ & $1$ & $4$ & $22$ & $52$ & $1.1825$ & $1855$ &  & $20$ & $57$ & $15$ & $0.8731$\\[3pt]
$1816$ & $0$ & $10$ & $11$ & $43$ & $0.4248$ & $1856$ &  & $2$ & $46$ & $7$ & $0.1154$\\
$1817$ & $0$ & $16$ & $0$ & $35$ & $0.6671$ & $1857$ &  & $8$ & $34$ & $58$ & $0.3576$\\
$1818$ & $0$ & $21$ & $49$ & $27$ & $0.9093$ & $1858$ &  & $14$ & $23$ & $50$ & $0.5999$\\
$1819$ & $1$ & $3$ & $38$ & $18$ & $1.1516$ & $1859$ &  & $20$ & $12$ & $41$ & $0.8422$\\
$1820$ & $0$ & $9$ & $27$ & $10$ & $0.3939$ & $1860$ &  & $2$ & $1$ & $33$ & $0.0844$\\[3pt]
$1821$ & $0$ & $15$ & $16$ & $1$ & $0.6361$ & $1864$ &  & $1$ & $16$ & $59$ & $0.0535$\\
$1822$ & $0$ & $21$ & $4$ & $53$ & $0.8784$ & $1868$ & $+0$ & $0$ & $32$ & $26$ & $+0.0225$\\
$1823$ & $1$ & $2$ & $53$ & $44$ & $1.1207$ & $1872$ & $-0$ & $0$ & $12$ & $8$ & $-0.0084$\\
$1824$ & $0$ & $8$ & $42$ & $36$ & $0.3629$ & $1876$ & $+0$ & $0$ & $56$ & $42$ & $+0.0393$\\
$1825$ & $0$ & $14$ & $31$ & $28$ & $0.6052$ & $1880$ &  & $1$ & $41$ & $15$ & $0.0703$\\[3pt]
$1826$ & $0$ & $20$ & $20$ & $19$ & $0.8474$ & $1884$ &  & $2$ & $25$ & $49$ & $0.1012$\\
$1827$ & $1$ & $2$ & $9$ & $11$ & $1.0897$ & $1888$ &  & $3$ & $10$ & $23$ & $0.1322$\\
$1828$ & $0$ & $7$ & $58$ & $2$ & $0.3320$ & $1892$ &  &  & $54$ & $57$ & $0.1631$\\
$1829$ & $0$ & $13$ & $46$ & $54$ & $0.5742$ & $1896$ &  & $4$ & $39$ & $30$ & $0.1941$\\
$1830$ & $0$ & $19$ & $35$ & $46$ & $0.8165$ & $1900$ & $+0$ & $18$ & $35$ & $56$ & $+0.7750$\\
\end{tabular}
\endgroup
\end{center}
\clearpage
\SourceLeaf{182}%
\begin{center}
II. Tafel.

\medskip\emph{Die Sonnen-Coordinaten.}\medskip

\begingroup\fontsize{9}{10.2}\selectfont\setlength{\tabcolsep}{1.5pt}
\begin{tabular}{rr|r|r|r||rr|r|r|r||r|r|r|r||r|r|r|r}
\multicolumn{5}{c||}{Januar.}&\multicolumn{5}{c||}{Februar.}&\multicolumn{4}{c||}{März.}&\multicolumn{4}{c}{April.}\\[3pt]
\rotatebox{90}{Gemein Jahr}&\rotatebox{90}{Schalt-Jahr}&$x$&$y$&$z$&\rotatebox{90}{Gemein Jahr}&\rotatebox{90}{Schalt-Jahr}&$x$&$y$&$z$& &$x$&$y$&$z$& &$x$&$y$&$z$\\
&&$+$&$-$&$-$&&&$+$&$-$&$-$&&$+$&$\mp$&$\mp$&&$+$&$+$&$+$\\\hline
 & $0$ & $1614$ & $8897$ & $3862$ & $0$ & $1$ & $6598$ & $6718$ & $2916$ & $1$ & $9392$ & $2922$ & $1268$ & $1$ & $9783$ & $1918$ & $832$\\
$0$ & $1$ & $1786$ & $8869$ & $3850$ & $1$ & $2$ & $6727$ & $6611$ & $2869$ & $2$ & $9449$ & $2772$ & $1203$ & $2$ & $9748$ & $2072$ & $899$\\
$1$ & $2$ & $1958$ & $8839$ & $3836$ & $2$ & $3$ & $6855$ & $6502$ & $2822$ & $3$ & $9503$ & $2621$ & $1137$ & $3$ & $9711$ & $2226$ & $966$\\
$2$ & $3$ & $2129$ & $8805$ & $3822$ & $3$ & $4$ & $6980$ & $6391$ & $2774$ & $4$ & $9553$ & $2469$ & $1072$ & $4$ & $9671$ & $2379$ & $1033$\\
$3$ & $4$ & $2300$ & $8769$ & $3806$ & $4$ & $5$ & $7104$ & $6277$ & $2725$ & $5$ & $9602$ & $2316$ & $1005$ & $5$ & $9627$ & $2532$ & $1099$\\[3pt]
$4$ & $5$ & $2469$ & $8731$ & $3789$ & $5$ & $6$ & $7225$ & $6162$ & $2675$ & $6$ & $9647$ & $2163$ & $939$ & $6$ & $9581$ & $2684$ & $1165$\\
$5$ & $6$ & $2638$ & $8689$ & $3771$ & $6$ & $7$ & $7344$ & $6045$ & $2624$ & $7$ & $9689$ & $2008$ & $872$ & $7$ & $9533$ & $2835$ & $1230$\\
$6$ & $7$ & $2806$ & $8645$ & $3752$ & $7$ & $8$ & $7460$ & $5927$ & $2572$ & $8$ & $9728$ & $1854$ & $805$ & $8$ & $9481$ & $2985$ & $1296$\\
$7$ & $8$ & $2974$ & $8598$ & $3732$ & $8$ & $9$ & $7574$ & $5806$ & $2520$ & $9$ & $9765$ & $1699$ & $737$ & $9$ & $9427$ & $3134$ & $1360$\\
$8$ & $9$ & $3140$ & $8549$ & $3711$ & $9$ & $10$ & $7686$ & $5683$ & $2467$ & $10$ & $9798$ & $1543$ & $670$ & $10$ & $9369$ & $3283$ & $1425$\\[3pt]
$9$ & $10$ & $3306$ & $8497$ & $3688$ & $10$ & $11$ & $7796$ & $5559$ & $2413$ & $11$ & $9828$ & $1387$ & $602$ & $11$ & $9310$ & $3430$ & $1489$\\
$10$ & $11$ & $3470$ & $8442$ & $3664$ & $11$ & $12$ & $7903$ & $5433$ & $2358$ & $12$ & $9856$ & $1230$ & $534$ & $12$ & $9247$ & $3576$ & $1552$\\
$11$ & $12$ & $3633$ & $8384$ & $3639$ & $12$ & $13$ & $8008$ & $5306$ & $2303$ & $13$ & $9881$ & $1073$ & $466$ & $13$ & $9182$ & $3722$ & $1615$\\
$12$ & $13$ & $3795$ & $8324$ & $3613$ & $13$ & $14$ & $8111$ & $5177$ & $2247$ & $14$ & $9902$ & $916$ & $397$ & $14$ & $9114$ & $3866$ & $1678$\\
$13$ & $14$ & $3956$ & $8262$ & $3586$ & $14$ & $15$ & $8210$ & $5046$ & $2190$ & $15$ & $9921$ & $758$ & $329$ & $15$ & $9043$ & $4009$ & $1740$\\[3pt]
$14$ & $15$ & $4116$ & $8196$ & $3558$ & $15$ & $16$ & $8308$ & $4913$ & $2133$ & $16$ & $9936$ & $600$ & $260$ & $16$ & $8969$ & $4151$ & $1802$\\
$15$ & $16$ & $4275$ & $8129$ & $3528$ & $16$ & $17$ & $8403$ & $4780$ & $2075$ & $17$ & $9949$ & $442$ & $192$ & $17$ & $8894$ & $4292$ & $1863$\\
$16$ & $17$ & $4432$ & $8058$ & $3498$ & $17$ & $18$ & $8495$ & $4644$ & $2016$ & $18$ & $9959$ & $284$ & $123$ & $18$ & $8815$ & $4431$ & $1923$\\
$17$ & $18$ & $4587$ & $7986$ & $3466$ & $18$ & $19$ & $8584$ & $4508$ & $1956$ & $19$ & $9965$ & $-126$ & $-55$ & $19$ & $8734$ & $4569$ & $1983$\\
$18$ & $19$ & $4742$ & $7910$ & $3433$ & $19$ & $20$ & $8671$ & $4369$ & $1897$ & $20$ & $9969$ & $+33$ & $+14$ & $20$ & $8650$ & $4706$ & $2043$\\
 &  &  &  &  &  &  &  &  & $1836$ &  &  &  &  &  &  &  & \\
$19$ & $20$ & $4895$ & $7833$ & $3400$ & $20$ & $21$ & $8756$ & $4230$ &  & $21$ & $9970$ & $191$ & $83$ & $21$ & $8564$ & $4842$ & $2102$\\
$20$ & $21$ & $5046$ & $7752$ & $3365$ & $21$ & $22$ & $8837$ & $4089$ & $1775$ & $22$ & $9968$ & $349$ & $152$ & $22$ & $8475$ & $4976$ & $2160$\\
$21$ & $22$ & $5196$ & $7670$ & $3329$ & $22$ & $23$ & $8916$ & $3948$ & $1713$ & $23$ & $9962$ & $507$ & $220$ & $23$ & $8384$ & $5109$ & $2217$\\
$22$ & $23$ & $5344$ & $7585$ & $3292$ & $23$ & $24$ & $8993$ & $3804$ & $1651$ & $24$ & $9954$ & $665$ & $289$ & $24$ & $8291$ & $5240$ & $2274$\\
$23$ & $24$ & $5490$ & $7498$ & $3254$ & $24$ & $25$ & $9066$ & $3660$ & $1588$ & $25$ & $9943$ & $823$ & $357$ & $25$ & $8195$ & $5369$ & $2331$\\[3pt]
$24$ & $25$ & $5635$ & $7408$ & $3215$ & $25$ & $26$ & $9137$ & $3514$ & $1525$ & $26$ & $9929$ & $980$ & $426$ & $26$ & $8096$ & $5498$ & $2386$\\
$25$ & $26$ & $5778$ & $7316$ & $3175$ & $26$ & $27$ & $9205$ & $3368$ & $1462$ & $27$ & $9912$ & $1138$ & $494$ & $27$ & $7996$ & $5624$ & $2441$\\
$26$ & $27$ & $5920$ & $7222$ & $3035$ & $27$ & $28$ & $9270$ & $3220$ & $1398$ & $28$ & $9892$ & $1295$ & $562$ & $28$ & $7893$ & $5749$ & $2495$\\
$27$ & $28$ & $6059$ & $7125$ & $3093$ & $28$ & $29$ & $9332$ & $3072$ & $1333$ & $29$ & $9869$ & $1451$ & $630$ & $29$ & $7787$ & $5872$ & $2549$\\
$28$ & $29$ & $6197$ & $7027$ & $3050$ & $29$ & $30$ & $9392$ & $2922$ & $1268$ & $30$ & $9843$ & $1607$ & $698$ & $30$ & $7680$ & $5994$ & $2601$\\[3pt]
$29$ & $30$ & $6332$ & $6926$ & $3006$ &  &  &  &  &  & $31$ & $9814$ & $1763$ & $765$ & $31$ & $7570$ & $6113$ & $2653$\\
$30$ & $31$ & $6466$ & $6823$ & $2962$ &  &  &  &  &  & $32$ & $9783$ & $1918$ & $832$ &  &  &  & \\
$31$ & $32$ & $6598$ & $6718$ & $2916$ &  &  &  &  &  &  &  &  &  &  &  &  & \\
\end{tabular}\endgroup
\end{center}
\clearpage
\SourceLeaf{183}%
\begin{center}
II. Tafel.

\medskip\emph{Die Sonnen-Coordinaten.}\medskip

\begingroup\fontsize{9}{10.2}\selectfont\setlength{\tabcolsep}{2pt}
\begin{tabular}{r|r|r|r||r|r|r|r||r|r|r|r||r|r|r|r}
\multicolumn{4}{c||}{Mai.} & \multicolumn{4}{c||}{Juni.} & \multicolumn{4}{c||}{Juli.} & \multicolumn{4}{c}{August.}\\[3pt]
 & $x$ & $y$ & $z$ &  & $x$ & $y$ & $z$ &  & $x$ & $y$ & $z$ &  & $x$ & $y$ & $z$\\
 & $+$ & $+$ & $+$ &  & $\pm$ & $+$ & $+$ &  & $-$ & $+$ & $+$ &  & $-$ & $+$ & $+$\\\hline
$1$ & $7570$ & $6113$ & $2653$ & $1$ & $3271$ & $8811$ & $3824$ & $1$ & $1737$ & $9190$ & $3989$ & $1$ & $6443$ & $7188$ & $3120$\\
$2$ & $7458$ & $6231$ & $2705$ & $2$ & $3110$ & $8861$ & $3846$ & $2$ & $1903$ & $9162$ & $3977$ & $2$ & $6573$ & $7088$ & $3076$\\
$3$ & $7344$ & $6348$ & $2755$ & $3$ & $2949$ & $8909$ & $3867$ & $3$ & $2069$ & $9132$ & $3964$ & $3$ & $6700$ & $6985$ & $3032$\\
$4$ & $7228$ & $6462$ & $2805$ & $4$ & $2787$ & $8954$ & $3886$ & $4$ & $2234$ & $9099$ & $3949$ & $4$ & $6825$ & $6880$ & $2986$\\
$5$ & $7110$ & $6575$ & $2854$ & $5$ & $2624$ & $8996$ & $3905$ & $5$ & $2399$ & $9063$ & $3934$ & $5$ & $6949$ & $6773$ & $2940$\\[3pt]
$6$ & $6990$ & $6685$ & $2902$ & $6$ & $2460$ & $9036$ & $3922$ & $6$ & $2563$ & $9025$ & $3917$ & $6$ & $7070$ & $6665$ & $2893$\\
$7$ & $6867$ & $6794$ & $2949$ & $7$ & $2295$ & $9074$ & $3938$ & $7$ & $2726$ & $8985$ & $3900$ & $7$ & $7189$ & $6554$ & $2845$\\
$8$ & $6743$ & $6901$ & $2995$ & $8$ & $2130$ & $9108$ & $3953$ & $8$ & $2889$ & $8942$ & $3881$ & $8$ & $7307$ & $6442$ & $2796$\\
$9$ & $6616$ & $7006$ & $3041$ & $9$ & $1965$ & $9141$ & $3967$ & $9$ & $3050$ & $8896$ & $3861$ & $9$ & $7422$ & $6327$ & $2746$\\
$10$ & $6488$ & $7108$ & $3085$ & $10$ & $1798$ & $9170$ & $3980$ & $10$ & $3211$ & $8848$ & $3840$ & $10$ & $7535$ & $6211$ & $2696$\\[3pt]
$11$ & $6358$ & $7209$ & $3129$ & $11$ & $1632$ & $9197$ & $3992$ & $11$ & $3371$ & $8797$ & $3818$ & $11$ & $7647$ & $6094$ & $2645$\\
$12$ & $6226$ & $7308$ & $3172$ & $12$ & $1464$ & $9222$ & $4003$ & $12$ & $3530$ & $8744$ & $3795$ & $12$ & $7755$ & $5974$ & $2593$\\
$13$ & $6093$ & $7404$ & $3214$ & $13$ & $1297$ & $9244$ & $4012$ & $13$ & $3688$ & $8689$ & $3771$ & $13$ & $7862$ & $5853$ & $2540$\\
$14$ & $5958$ & $7499$ & $3255$ & $14$ & $1129$ & $9263$ & $4021$ & $14$ & $3845$ & $8631$ & $3746$ & $14$ & $7967$ & $5730$ & $2487$\\
$15$ & $5821$ & $7591$ & $3295$ & $15$ & $961$ & $9280$ & $4028$ & $15$ & $4001$ & $8570$ & $3720$ & $15$ & $8069$ & $5605$ & $2433$\\[3pt]
$16$ & $5682$ & $7681$ & $3334$ & $16$ & $792$ & $9294$ & $4034$ & $16$ & $4156$ & $8507$ & $3692$ & $16$ & $8169$ & $5479$ & $2378$\\
$17$ & $5541$ & $7769$ & $3372$ & $17$ & $623$ & $9305$ & $4039$ & $17$ & $4309$ & $8442$ & $3664$ & $17$ & $8266$ & $5351$ & $2323$\\
$18$ & $5400$ & $7855$ & $3409$ & $18$ & $454$ & $9314$ & $4043$ & $18$ & $4462$ & $8374$ & $3635$ & $18$ & $8361$ & $5222$ & $2267$\\
$19$ & $5256$ & $7939$ & $3446$ & $19$ & $285$ & $9320$ & $4045$ & $19$ & $4613$ & $8304$ & $3604$ & $19$ & $8454$ & $5091$ & $2210$\\
$20$ & $5111$ & $8020$ & $3481$ & $20$ & $+116$ & $9324$ & $4047$ & $20$ & $4762$ & $8232$ & $3573$ & $20$ & $8544$ & $4959$ & $2153$\\[3pt]
$21$ & $4965$ & $8099$ & $3515$ & $21$ & $-53$ & $9325$ & $4047$ & $21$ & $4911$ & $8157$ & $3541$ & $21$ & $8632$ & $4826$ & $2095$\\
$22$ & $4817$ & $8175$ & $3548$ & $22$ & $222$ & $9323$ & $4047$ & $22$ & $5058$ & $8080$ & $3507$ & $22$ & $8718$ & $4691$ & $2036$\\
$23$ & $4668$ & $8250$ & $3581$ & $23$ & $391$ & $9319$ & $4045$ & $23$ & $5204$ & $8001$ & $3473$ & $23$ & $8801$ & $4554$ & $1977$\\
$24$ & $4517$ & $8322$ & $3612$ & $24$ & $560$ & $9312$ & $4042$ & $24$ & $5348$ & $7919$ & $3437$ & $24$ & $8881$ & $4417$ & $1917$\\
$25$ & $4365$ & $8391$ & $3642$ & $25$ & $729$ & $9302$ & $4038$ & $25$ & $5490$ & $7836$ & $3401$ & $25$ & $8959$ & $4278$ & $1857$\\[3pt]
$26$ & $4212$ & $8459$ & $3671$ & $26$ & $898$ & $9290$ & $4032$ & $26$ & $5632$ & $7750$ & $3364$ & $26$ & $9034$ & $4137$ & $1796$\\
$27$ & $4058$ & $8523$ & $3700$ & $27$ & $1066$ & $9275$ & $4026$ & $27$ & $5771$ & $7661$ & $3325$ & $27$ & $9107$ & $3996$ & $1735$\\
$28$ & $3903$ & $8586$ & $3727$ & $28$ & $1235$ & $9258$ & $4018$ & $28$ & $5909$ & $7571$ & $3286$ & $28$ & $9177$ & $3854$ & $1673$\\
$29$ & $3746$ & $8646$ & $3753$ & $29$ & $1402$ & $9238$ & $4010$ & $29$ & $6045$ & $7479$ & $3246$ & $29$ & $9245$ & $3710$ & $1610$\\
$30$ & $3589$ & $8703$ & $3778$ & $30$ & $1570$ & $9215$ & $4000$ & $30$ & $6180$ & $7384$ & $3205$ & $30$ & $9309$ & $3565$ & $1548$\\[3pt]
$31$ & $3430$ & $8758$ & $3802$ & $31$ & $1737$ & $9190$ & $3989$ & $31$ & $6312$ & $7287$ & $3163$ & $31$ & $9371$ & $3420$ & $1484$\\
$32$ & $3271$ & $8811$ & $3824$ &  &  &  &  & $32$ & $6443$ & $7188$ & $3120$ & $32$ & $9431$ & $3273$ & $1421$\\
\end{tabular}\endgroup
\end{center}
\clearpage
\SourceLeaf{184}%
\begin{center}
II. Tafel.

\medskip\emph{Die Sonnen-Coordinaten.}\medskip

\begingroup\fontsize{9}{10.2}\selectfont\setlength{\tabcolsep}{2pt}
\begin{tabular}{r|r|r|r||r|r|r|r||r|r|r|r||r|r|r|r}
\multicolumn{4}{c||}{September.} & \multicolumn{4}{c||}{October.} & \multicolumn{4}{c||}{November.} & \multicolumn{4}{c}{December.}\\[3pt]
 & $x$ & $y$ & $z$ &  & $x$ & $y$ & $z$ &  & $x$ & $y$ & $z$ &  & $x$ & $y$ & $z$\\
 & $-$ & $\pm$ & $\pm$ &  & $-$ & $-$ & $-$ &  & $-$ & $-$ & $-$ &  & $\mp$ & $-$ & $-$\\\hline
$1$ & $9431$ & $3273$ & $1421$ & $1$ & $9889$ & $1370$ & $595$ & $1$ & $7661$ & $5774$ & $2506$ & $1$ & $3422$ & $8477$ & $3679$\\
$2$ & $9488$ & $3125$ & $1357$ & $2$ & $9859$ & $1526$ & $662$ & $2$ & $7548$ & $5894$ & $2558$ & $2$ & $3257$ & $8530$ & $3702$\\
$3$ & $9541$ & $2977$ & $1292$ & $3$ & $9826$ & $1681$ & $729$ & $3$ & $7432$ & $6013$ & $2610$ & $3$ & $3091$ & $8580$ & $3724$\\
$4$ & $9593$ & $2827$ & $1227$ & $4$ & $9790$ & $1835$ & $796$ & $4$ & $7315$ & $6130$ & $2661$ & $4$ & $2925$ & $8628$ & $3745$\\
$5$ & $9641$ & $2677$ & $1162$ & $5$ & $9751$ & $1989$ & $863$ & $5$ & $7195$ & $6245$ & $2711$ & $5$ & $2757$ & $8673$ & $3765$\\[3pt]
$6$ & $9687$ & $2526$ & $1097$ & $6$ & $9710$ & $2142$ & $930$ & $6$ & $7072$ & $6358$ & $2760$ & $6$ & $2589$ & $8716$ & $3783$\\
$7$ & $9729$ & $2375$ & $1031$ & $7$ & $9665$ & $2295$ & $996$ & $7$ & $6948$ & $6470$ & $2808$ & $7$ & $2419$ & $8755$ & $3800$\\
$8$ & $9769$ & $2222$ & $965$ & $8$ & $9618$ & $2447$ & $1062$ & $8$ & $6822$ & $6579$ & $2856$ & $8$ & $2249$ & $8792$ & $3816$\\
$9$ & $9807$ & $2069$ & $898$ & $9$ & $9567$ & $2598$ & $1128$ & $9$ & $6693$ & $6686$ & $2902$ & $9$ & $2079$ & $8826$ & $3831$\\
$10$ & $9841$ & $1915$ & $831$ & $10$ & $9514$ & $2749$ & $1193$ & $10$ & $6562$ & $6792$ & $2948$ & $10$ & $1907$ & $8858$ & $3845$\\[3pt]
$11$ & $9872$ & $1761$ & $764$ & $11$ & $9458$ & $2898$ & $1258$ & $11$ & $6430$ & $6895$ & $2993$ & $11$ & $1735$ & $8887$ & $3857$\\
$12$ & $9901$ & $1606$ & $697$ & $12$ & $9400$ & $3047$ & $1323$ & $12$ & $6295$ & $6996$ & $3037$ & $12$ & $1563$ & $8913$ & $3869$\\
$13$ & $9927$ & $1451$ & $630$ & $13$ & $9338$ & $3195$ & $1387$ & $13$ & $6159$ & $7095$ & $3079$ & $13$ & $1390$ & $8936$ & $3879$\\
$14$ & $9949$ & $1295$ & $562$ & $14$ & $9274$ & $3342$ & $1451$ & $14$ & $6020$ & $7192$ & $3121$ & $14$ & $1217$ & $8956$ & $3887$\\
$15$ & $9969$ & $1140$ & $495$ & $15$ & $9206$ & $3488$ & $1514$ & $15$ & $5880$ & $7286$ & $3163$ & $15$ & $1043$ & $8974$ & $3895$\\[3pt]
$16$ & $9986$ & $983$ & $427$ & $16$ & $9137$ & $3633$ & $1577$ & $16$ & $5738$ & $7379$ & $3203$ & $16$ & $869$ & $8989$ & $3902$\\
$17$ & $10000$ & $827$ & $359$ & $17$ & $9064$ & $3777$ & $1639$ & $17$ & $5594$ & $7469$ & $3242$ & $17$ & $694$ & $9001$ & $3907$\\
$18$ & $10011$ & $670$ & $291$ & $18$ & $8989$ & $3919$ & $1701$ & $18$ & $5449$ & $7556$ & $3280$ & $18$ & $520$ & $9010$ & $3911$\\
$19$ & $10020$ & $513$ & $223$ & $19$ & $8911$ & $4061$ & $1763$ & $19$ & $5302$ & $7642$ & $3317$ & $19$ & $345$ & $9017$ & $3914$\\
$20$ & $10025$ & $356$ & $154$ & $20$ & $8830$ & $4201$ & $1823$ & $20$ & $5153$ & $7725$ & $3353$ & $20$ & $-170$ & $9020$ & $3915$\\[3pt]
$21$ & $10027$ & $198$ & $86$ & $21$ & $8746$ & $4340$ & $1884$ & $21$ & $5002$ & $7806$ & $3388$ & $21$ & $+5$ & $9021$ & $3916$\\
$22$ & $10027$ & $+41$ & $+18$ & $22$ & $8660$ & $4478$ & $1944$ & $22$ & $4851$ & $7884$ & $3422$ & $22$ & $180$ & $9019$ & $3915$\\
$23$ & $10023$ & $-116$ & $-50$ & $23$ & $8572$ & $4614$ & $2003$ & $23$ & $4697$ & $7960$ & $3455$ & $23$ & $354$ & $9015$ & $3913$\\
$24$ & $10017$ & $274$ & $119$ & $24$ & $8481$ & $4749$ & $2061$ & $24$ & $4542$ & $8033$ & $3487$ & $24$ & $529$ & $9007$ & $3910$\\
$25$ & $10007$ & $431$ & $187$ & $25$ & $8387$ & $4883$ & $2119$ & $25$ & $4386$ & $8104$ & $3518$ & $25$ & $704$ & $8997$ & $3905$\\[3pt]
$26$ & $9995$ & $588$ & $255$ & $26$ & $8290$ & $5015$ & $2177$ & $26$ & $4228$ & $8173$ & $3547$ & $26$ & $878$ & $8984$ & $3899$\\
$27$ & $9979$ & $745$ & $323$ & $27$ & $8192$ & $5145$ & $2233$ & $27$ & $4069$ & $8239$ & $3576$ & $27$ & $1052$ & $8968$ & $3892$\\
$28$ & $9961$ & $902$ & $391$ & $28$ & $8090$ & $5274$ & $2289$ & $28$ & $3909$ & $8302$ & $3604$ & $28$ & $1226$ & $8949$ & $3884$\\
$29$ & $9940$ & $1058$ & $459$ & $29$ & $7987$ & $5402$ & $2344$ & $29$ & $3748$ & $8363$ & $3630$ & $29$ & $1399$ & $8927$ & $3875$\\
$30$ & $9916$ & $1214$ & $527$ & $30$ & $7881$ & $5527$ & $2399$ & $30$ & $3585$ & $8421$ & $3655$ & $30$ & $1572$ & $8903$ & $3864$\\[3pt]
$31$ & $9889$ & $1370$ & $595$ & $31$ & $7772$ & $5651$ & $2453$ & $31$ & $3422$ & $8477$ & $3679$ & $31$ & $1744$ & $8876$ & $3853$\\
 &  &  &  & $32$ & $7661$ & $5774$ & $2506$ &  &  &  &  & $32$ & $1916$ & $8846$ & $3840$\\
\end{tabular}\endgroup
\end{center}
\clearpage
\SourceLeaf{185}%
\begin{center}
III. Tafel.

\medskip\emph{Die Säcular-Änderungen der Sonnen-Coordinaten.}\medskip

\begin{tabular}{lr|r|r|r}
&&$x$&$y$&$z$\\\hline
Januar &$0$&$-10$&$-3$&$+1$\\
&$20,\ 19$&$9$&$4$&$+1$\\
Februar &$9,\ 8$&$8$&$5$&$0$\\
&$29,\ 28$&$6$&$4$&$-1$\\
März &$20$&$5$&$2$&$1$\\
April &$9$&$5$&$-1$&$1$\\
&$29$&$6$&$+1$&$1$\\
Mai &$19$&$8$&$1$&$2$\\
Junius &$8$&$9$&$+1$&$2$\\
&$28$&$10$&$-1$&$3$\\
Julius &$18$&$9$&$2$&$3$\\
August &$7$&$8$&$3$&$3$\\
&$27$&$6$&$3$&$3$\\
September &$16$&$5$&$3$&$-1$\\
October &$6$&$5$&$1$&$0$\\
&$26$&$6$&$0$&$+1$\\
November &$15$&$8$&$0$&$2$\\
December &$5$&$9$&$1$&$2$\\
&$25$&$10$&$2$&$2$
\end{tabular}
\end{center}
